\section{Líneas de trabajo futuro}

El desarrollo de Voctomix 2.0 ha puesto de manifiesto varias áreas de mejora e
investigación que quedan fuera del alcance de este Trabajo de Fin de Grado pero que
constituyen líneas de trabajo futuro de interés:

\begin{itemize}
  \item \textbf{Orquestación con Kubernetes:} migrar el fichero \texttt{docker-compose.yml}
        a un \textit{chart} Helm permitiría escalar voctocore horizontalmente para
        producciones multi-sala, donde varias salas de conferencias comparten la
        misma infraestructura de streaming.

  \item \textbf{Interfaz web de control:} sustituir voctogui por una aplicación web
        de página única (SPA) basada en React o Vue que consuma la API REST del notario
        eliminaría la dependencia de GTK y permitiría controlar el sistema desde
        cualquier dispositivo con navegador, incluidos tabletas y teléfonos móviles.

  \item \textbf{Detección automática de silencio:} integrar un clasificador de audio
        que monitorice continuamente los niveles de entrada y alerte al operador cuando
        una fuente pierde señal (micrófono desconectado, cable en mal estado), reduciendo
        los incidentes técnicos durante producciones en directo.

  \item \textbf{Grabación por segmentos:} implementar un \textit{sink} de grabación que
        genere ficheros MKV por ponencia de forma automática, sincronizado con los
        eventos del notario (inicio/fin de ponencia), eliminando la necesidad de post-proceso
        manual de los ficheros de grabación.

  \item \textbf{Soporte NDI:} incorporar el protocolo NDI (\textit{Network Device Interface})
        como fuente de entrada alternativa a TCP/MKV, facilitando la integración con
        cámaras de red y \textit{hardware} de producción moderno sin necesidad de
        adaptadores de señal adicionales.
\end{itemize}
